\section{Cómo hacerlo crecer}\label{sec:crecer}

KalaScript no está terminado ni debe estarlo: tiene que crecer con mis apuntes
reales. Esta sección fija cómo se añade, cambia o retira un signo sin que los
apuntes viejos dejen de leerse.

\subsection{Leyenda personal}\label{sec:leyenda}
La primera página de cada cuaderno es la leyenda: las abreviaturas propias, los
nombres frecuentes del curso y los signos candidatos que estoy probando. Si
alguien (o yo dentro de un año) abre el cuaderno, ahí está todo lo que no
aparece en este manual.

\subsection{Cuándo merece un signo nuevo}
Durante la clase no invento nada: escribo la palabra y la subrayo con una línea
ondulada. Al repasar, cuento cuántas veces la subrayé. Una idea pasa a
candidata si cumple las cuatro condiciones:
\begin{enumerate}
\item \textbf{Frecuencia.} Aparece varias veces por semana, en más de un curso.
\item \textbf{Ganancia.} El signo es claramente más rápido o más preciso que
 la palabra o su abreviatura.
\item \textbf{Hueco.} Ningún signo, símbolo matemático o abreviatura existente
 lo cubre ya.
\item \textbf{Distinción.} Su forma no se confunde, escrita rápido, con otro signo
 ni con una letra.
\end{enumerate}

\subsection{Cómo se prueba}
\begin{enumerate}
\item Dibujo dos o tres formas posibles y elijo la que se escriba en menos trazos.
\item La anoto en la leyenda personal con un código provisional (\texttt{C1},
 \texttt{C2}\ldots) y la uso dos semanas.
\item Al final, releo los apuntes de esas dos semanas sin la leyenda. Si la
 entiendo siempre y no la confundí, la registro con el siguiente código K libre.
 Si no, la descarto.
\end{enumerate}

\subsection{Reglas de estabilidad}
\begin{itemize}
\item Un código K nunca se reutiliza para otro significado.
\item Cambiar el dibujo de un signo está permitido si el significado no cambia;
 el dibujo anterior se conserva en el anexo para leer apuntes viejos.
\item Cambiar el significado exige retirar el signo y crear uno nuevo, salvo que
 el cambio solo amplíe su uso y quede escrito en el anexo de cambios.
\item Un signo que no uso en todo un ciclo es candidato a retirarse.
\end{itemize}

\subsection{Qué medir}
Para saber si KalaScript funciona de verdad, el cuaderno de práctica incluye
dos pruebas que conviene repetir cada mes:
\begin{itemize}
\item \textbf{Velocidad.} Copio el mismo párrafo en español normal y en
 KalaScript, con cronómetro.
\item \textbf{Relectura.} Leo apuntes de hace dos semanas sin el manual y anoto
 cada signo o línea que no entendí a la primera.
\end{itemize}
Si un signo aparece una y otra vez en la lista de relectura, el problema es el
signo, no mi memoria.

\subsection{Ideas para más adelante}
\begin{itemize}
\item Variantes cursivas que unan un conector con la palabra siguiente sin
 levantar el lápiz.
\item Signos propios por materia (por ejemplo, para química o para música) en
 una hoja aparte, sin cargar el inventario común.
\item Una fuente tipográfica para escribir KalaScript en el computador.
\end{itemize}
